% Part: incompleteness
% Chapter: incompleteness-provability
% Section: godels-paper

\documentclass[../../../include/open-logic-section]{subfiles}

\begin{document}

\olfileid[pt]{inc}{inp}{gop}

\olsection{Comparação com o artigo original de G\"odel}

Vale a pena dedicar algum tempo ao artigo de G\"odel de 1931. A introdução
esboça as ideias que acabamos de discutir. Mesmo que você apenas percorra o
artigo rapidamente, é fácil perceber o que acontece em cada etapa: primeiro
G\"odel descreve o sistema formal $P$ (sintaxe, axiomas, regras de prova);
depois define as funções e relações recursivas primitivas; em seguida mostra
que $x B y$ é recursiva primitiva e argumenta que as funções e relações
recursivas primitivas são representadas em $\Th{P}$. Ele então passa a
provar o teorema da incompletude, como acima. Na Seção 3, mostra que podemos
tomar a asserção indemonstrável como uma sentença na linguagem da
aritmética. Essa é a origem do lema~$\beta$, que também usamos para lidar
com sequências ao mostrar que as funções recursivas são representáveis em
$\Th{Q}$. G\"odel não chega a isolar um conjunto mínimo de axiomas
suficiente, mas hoje sabemos que $\Th{Q}$ cumpre essa função. Finalmente,
na Seção 4, ele esboça uma prova do segundo teorema da incompletude.

\end{document}

